\begin{textsample}{POS Dim 5 – sc – Score 6.00 – sc/sc\_em\_55.txt}  \label{ex:f5_pos_039}
2.3.4 Orientações metodológicas As \textbf{sequências} didáticas permeiam a problematização e a contextualização, a pesquisa, a resolução de problemas mediante processos de observação, comparação, investigação, análise, interpretação, reflexão e síntese ( PCSC, 2014 ).
Sugere -se, para tanto : - criar núcleos de estudos e observatórios sobre as formas de violência contra as mulheres nos lugares de vivência no território catarinense ( este núcleo teria por função desenvolver estudos e pesquisas sobre o tema, promovendo fóruns de debates, campanhas e \textbf{palestras} ) ; - identificar como, e se, os sujeitos das culturas escolares reconhecem as várias formas de violência contra as mulheres nos lugares de vivência no território catarinense ; - considerando as dimensões dos diversos tipos de violência ( física, psicológica, moral, patrimonial ), eleger uma delas e investigar como realizar o encaminhamento das diferentes situações das vítimas, sob a perspectiva de gênero, nos lugares de vivência ( escola, bairro, cidade ) ; - identificar a possível existência de sistemas de proteção contra violências das mulheres e quais possam ser acionados, nesses casos, para proteger as vítimas nos lugares de vivência ( bairros e cidade ), identificando pontos positivos e negativos para a garantia desse direito protetivo ; - propor políticas públicas, criar projetos ou programas que melhorem a situação de escuta do sofrimento e de acolhida da violência contra as mulheres, considerando as interseccionalidades de gênero, populações LGBTQIAP+, etnia/raça, de classe, geracional, entre outras ; - criar grupos de memória e de escuta das mulheres ( estudantes, jovens, idosas, adultas ) ; - promover ações que oportunizem \textbf{entrevistas} com mulheres mais velhas, e de diferentes gerações, para a percepção e constatação dos diferentes costumes e padrões de comportamento definidos para as mulheres, nas dimensões históricas e geográficas, sociológicas e filosóficas, nos lugares de vivência e promover rodas de conversas intergeracionais nas escolas ou em lugares de vivência. 2.3.5 Recursos Tempo de planejamento para os(as) professores(as) e para a execução da atividade.
O tempo de execução também pressupõe capacidade de ocupar diferentes espaços formativos ( fazer visitas a locais, reservar espaços, etc. ), tempo de planejamento com materiais apropriados, como \textbf{livros}, bibliografias atualizadas, funcionamento de internet, entre outros. \textbf{Sugestão} : saídas a campo, laboratórios, visitas guiadas, entre outras, o que requer : organização de transporte ; pesquisas bibliográficas ; preparação de \textbf{entrevistas} e sua realização ; viagens de estudo a órgãos do Poder Judiciário e de Serviço Social ; visitas à comunidades tradicionais ; análise dos espaços de memória ; pesquisas on-line ; documentários ; filmes ; artigos e outros ; acesso à internet, especialmente às redes sociais através do \textbf{celular} ( que geralmente não consome pacotes de dados ) ; referências bibliográficas ; \textbf{celulares} e notebooks ; parcerias com instituições ( universidades e laboratórios de pesquisa sobre violência contra as mulheres ) ; acessibilidade, visitas de estudos nas comunidades, formação de professores, criação de projetos pelos(as) estudantes ; laboratórios, salas-ambiente.

% matched lemmas: celular, entrevista, livro, palestra, sequência, sugestão
\end{textsample}
